\section{\label{sec:method}
Renormalization of EFT via on-shell methods}

In this section, we first review the on-shell method for computing anomalous dimensions introduced in Ref.~\cite{Caron-Huot:2016cwu}. Then, we discuss two independent ways to perform 
phase-space integrals both via angular variables~\cite{Zwiebel:2011bx} and through the use of Stokes' 
Theorem~\cite{Mastrolia:2009dr}. 
%

\subsection{The on-shell method}
%
We consider an effective Lagrangian of the type
\begin{equation}
\mathcal L_{\text{EFT}} = \sum_i 
\frac{c_i}{\Lambda^{[\mathcal O_i]-4}} \,\mathcal O_i \, , 
\end{equation}
%
where $\mathcal{O}_i$ are local gauge-invariant operators, $c_i$ are the corresponding Wilson coefficients, and $\Lambda$ refers to the UV cut-off scale of our EFT.

{\it Form factors} of the operators $\mathcal{O}_i$ are generically defined as
%
\begin{equation}
    F_i(\vec n;q)= \frac{1}{\Lambda^{[\mathcal O_i]-4}}\mel*{\vec n}{\mathcal O_i(q)}{0}\,,
\end{equation}
%
namely as the matrix element between an outgoing on-shell state $\bra*{\vec n}=\bra*{1^{h_1},\dotsc,n^{h_n}}$ and an operator $\mathcal O_i$ that injects an additional off-shell momentum $q$.
Within the dimensional regularization scheme, form factors depend on the renormalization scale $\mu$, and satisfy the Callan-Symanzik equation
%
\begin{equation}
\label{eq:CSEq}
    \left(\delta_{ij} \mu\frac{\partial}{\partial \mu}+ \frac{\partial \beta_i}{\partial c_j}-\delta_{ij}\gamma_{i,\txt{IR}} +\delta_{ij}\beta_g\frac{\partial}{\partial g}\right)F_i=   0\,,
\end{equation}
%
where $g$ collectively denotes the couplings related to the renormalizable operators of our Lagrangian, while $\gamma_{i,\txt{IR}}$ is the infrared anomalous dimension. 
The renormalization of the operator $\mathcal{O}_i$ 
is described by
%
\begin{equation}
\beta_i(\{c_k\}) \equiv \mu \dv[]{c_i}{\mu} \,, 
\end{equation}
%
where $c_i$ are the Wilson coefficients of the effective Lagrangian $\mathcal{L}_{\text{EFT}}$.

%
Exploiting the analyticity of form factors, unitarity, and the CPT theorem, it can be shown that an elegant relation exists linking the action of the dilatation operator ($D$)
%
to the action of the $S$-matrix ($S$) on form factors~\cite{Caron-Huot:2016cwu}:
%
\begin{equation}
\label{eq:FundamentalRelation}
    e^{-i \pi D}F_i^* = SF_i^*
\end{equation}
%
where $S=\mathbf{1} + i\mathcal{M}$ while $D = \sum_{i}p_i\cdot \partial/\partial p_i$
(the sum runs over all particles $i$). 

It is precisely the combination of Eqs.~\eqref{eq:FundamentalRelation} and \eqref{eq:CSEq} that allows one to directly link the renormalization group coefficients to the $S$-matrix.
%
In particular, at one-loop order, it has been found that
%
\begin{equation}
\label{eq:masterformula}
\left(
\frac{\partial \beta_i^{(1)}}{\partial c_j}
-\delta_{ij}\gamma_{i,\txt{IR}}^{(1)} 
+\delta_{ij}\beta_g^{(1)}\frac{\partial}{\partial g}
\right)
F_i^{(0)} = -\frac{1}{\pi}(\mathcal M F_j)^{(1)}\,, 
\end{equation}
%
where the right-hand side of Eq.~\eqref{eq:masterformula} corresponds to a sum over all one-loop two-particle unitarity cuts,
%
\begin{align}
\label{eq:MF1Loop}
    (\mathcal M F_j)^{(1)}(1,\dots,n) = 
   &\sum_{k=2}^n \sum_{\{x,y\}}\int d\text{LIPS}_2
   \nonumber\\ 
   &\times\sum_{h_1,h_2}\!\! F_j^{(0)}\!(x^{h_1},y^{h_2},k \!+\! 1,\dots,n) \mathcal M^{(0)}\!(1,\dots,k;x^{h_1},y^{h_2}\!)\, ,
\end{align}
where $\mathcal M(\vec n; \vec m)=\mel*{\vec n}{\mathcal M}{\vec m}$, and $d\text{LIPS}_2$ is the (two-particle) Lorentz invariant phase-space measure. 
The corresponding cut-integral can be evaluated either by angular integration~\cite{Zwiebel:2011bx} or via Stokes' theorem~\cite{Mastrolia:2009dr},
which differ for the parameterizations of the on-shell loop momentum across the cut and for the employed integration techniques, 
as we will show in the following.

Two observations are in order. 
Let us first remark that, at one-loop level, the $\beta$ function does not contribute to minimal form factors, because 
the latter are expectation values of purely local products of fields, and, by definition, they are independent of the renormalizable couplings of the theory, which we have collectively denoted by $g$.
The contribution from the $\beta$ function of renormalizable couplings is however unavoidable at higher perturbative orders.
%
Secondly, we stress that, owing to the adopted dimensional regularization scheme, this method is sensitive only to the difference between UV and IR divergences. 
Therefore, in order to disentangle the renormalization group equations for the UV divergent part, the IR contribution must be computed independently (see Appendix~\ref{app:IRanomalousdimensions} for more details). 

The on-shell method, so far discussed, is well suited to compute anomalous dimensions in massless theories, where all 
EFT operators have the same mass-dimension, as shown in the case of Standard Model EFT~\cite{EliasMiro:2020tdv}.  

In particular, in the case of linear operator mixing, 
Eq.~\eqref{eq:masterformula}
becomes
%
\begin{equation}
\label{eq:masterformula1L}
\left(
\gamma_{i \leftarrow j}^{(1)}
-\delta_{ij}\gamma_{i,\txt{IR}}^{(1)} 
\right)
F_i|_{*}^{(0)} = -\frac{1}{\pi}(\mathcal M F_j)|_{*}^{(1)}
\end{equation}
%
which has been evaluated at the Gaussian fixed point $(*)$, where all the Wilson coefficients $c_i$ are vanishing.
Moreover, $\gamma_{i \leftarrow j}^{(1)}$ is obtained from
the Taylor expansion of the renormalization group equations for the Wilson coefficients $c_i$
%
\begin{equation}
    \mu\dv[]{c_i}{\mu} = 
    \gamma_{i\leftarrow j}c_j +
    \frac{1}{2}\gamma_{i\leftarrow j,k}c_jc_k + \dots\,,
\end{equation}
%
where 
%
\begin{equation}
    \gamma_{i\leftarrow j_1, \dotsc, j_n} = \left.\frac{\partial^n \beta_i}{\partial c_{j_1}\dotsm \partial c_{j_n}}\right |_{*}\,.
\end{equation}
%

The result of Eq.~\eqref{eq:masterformula1L} can be easily applied to the case of non-linear mixing among operators with different dimensions, which is of interest to our study. At one-loop order, one can find the following expression~\cite{Bresciani:2023jsu} 
%
\begin{equation}
\label{eq:mastergammaijk}
\gamma_{i\leftarrow j,k}^{(1)}F_i|_{*}^{(0)} = -\frac{1}{\pi}\left.\frac{\partial}{\partial c_k}\right|_{c_k=0}(\mathcal M F_j)|_{*,c_k\neq 0}^{(1)} \,,
\end{equation}
%
where $j,k\neq i$ and
%
\begin{align}
    \gamma_{i\leftarrow j,k}&=\left. \frac{\partial^2 \beta_i}{\partial c_j \partial c_k}\right|_{*} = \left.\frac{\partial}{\partial c_k}\right|_{c_k=0} \left.\frac{\partial \beta_i}{\partial c_j}\right|_{*,c_k\neq 0} \,.
    \label{eq:gamma_ijk}
\end{align}
% 
Moreover, by making use of the Higgs low-energy theorem~\cite{Ellis:1975ap,Shifman:1979eb}, it is possible to include leading mass effects while still working in a massless formalism~\cite{Bresciani:2023jsu}. 
%
In practice, whenever an amplitude requires $N$ fermion mass insertions not to vanish, we consider an equivalent amplitude entailing $N$ extra massless Higgs fields, where 
%
\begin{equation}
N = 4-[\mathcal{O}_{i}] +\sum^n_{k=1} ([\mathcal{O}_{j_k}]-4)\geq 0
\label{eq:superficial_div}
\end{equation}
%
corresponds to the superficial degree of divergence associated with the loop diagram under consideration~\cite{Bresciani:2023jsu}.
%
Then, the anomalous dimension
$\gamma_{i\leftarrow j_1, \dotsc, j_n}$ 
is obtained by renormalizing the operator $(h/v)^N\mathcal{O}_i/N!$
instead of $\mathcal{O}_i$~\cite{Bresciani:2023jsu}.
Instead, for $N < 0$, $\gamma_{i\leftarrow j_1, \dotsc, j_n}$ does vanish.

%
The Lorentz-invariant phase space measure appearing in Eq.~\eqref{eq:MF1Loop} can be parameterized in different ways, depending on the employed method of integration. 
Here we will focus on two possible integration techniques: one based on an angular parameterization of the phase space and the other relying on the use of Stokes' theorem.
Whereas the former has the virtue of being quite intuitive, the latter turns out to be more suited to our purposes. Based on an elegant mathematical result, it offers a simpler and more direct way of carrying out phase-space integrations, as shown in several explicit examples.

\subsection{Phase-space integrals via angular variables}
\label{sec:phasespace_Angular}

The integration with angular variables relies on a parameterization of the virtual phase space
that is realized through the application of the following spinor rotation matrix~\cite{Zwiebel:2011bx}:
\begin{eqnarray}\label{eq:angular_param}
    \begin{pmatrix}
 \lambda_x \\
 \lambda_y
    \end{pmatrix}
     = 
    \begin{pmatrix}
 \cos\theta & -\sin\theta e^{i\phi} \\ 
 \sin\theta e^{-i\phi} & \cos\theta
\end{pmatrix}
\begin{pmatrix}
 \lambda_{a} \\
 \lambda_{b}
    \end{pmatrix}\,,
\end{eqnarray}
where $ \lambda_{a},\lambda_{b}$ correspond to the external momenta $p_a,p_b$,
and the integration measure is 
\begin{equation}
    \int d\text{LIPS}_2= \frac{1}{8 \pi}\int_0^{2 \pi} \frac{d\phi}{2\pi}\int_0^{\pi/2} 2 \cos\theta \sin\theta \, d\theta  \, . 
\end{equation}
%
An additional factor of $1/2$ has to be included in the measure if the two particles across the cut are indistinguishable.

%
Sometimes it is convenient to use the following parameterization
%
\begin{equation}
    \begin{pmatrix}
        \lambda_x \\ \lambda_y
    \end{pmatrix}
    =
    \frac{1}{\sqrt{1+t^2}}
    \begin{pmatrix}
        1 & -tz \\
        t/z & 1
    \end{pmatrix}
    \begin{pmatrix}
        \lambda_a \\ \lambda_b
    \end{pmatrix}\,,
\end{equation}
%
which rationalizes the integrand function and is derived from Eq.~\eqref{eq:angular_param} by setting $t=\tan\theta$ and $z=e^{i\phi}$.
The corresponding integration measure is given by
\begin{equation}
\label{eq:hybbrid}
    \int d\text{LIPS}_2 = \frac{1}{8\pi}\int_0^\infty \frac{2t \, dt}{(1+t^2)^2}\oint_{|z|=1}\frac{dz}{2\pi i z}\,.
\end{equation}


\subsection{Phase-space integrals via Stokes' Theorem}
\label{sec:phasespace_Stokes}
%
One-loop Feynman integrals, as well as scattering amplitudes in dimensional regularization, with $d= 4- \epsilon $ space-time dimensions,  can be decomposed in a finite basis of scalar integrals, known as master integrals.
Remarkably, up to order $\mathcal{O}(\epsilon^0)$, for any one-loop $n$-point amplitude, such master integrals are 4-point, 3-point, 2-point, and 1-point functions.
The latter do not contribute to processes with massless internal states; therefore, the UV singularities of massless 1-loop amplitudes,
together with collinear IR divergences,
are entirely contained in the 2-point function, and
they are proportional to the associated decomposition coefficient.
Singularities associated with 3- and 4-point functions are instead related only to IR divergences. 
Various techniques have been developed to evaluate the decomposition coefficients, including integration-by-parts identities \cite{Tkachov:1981wb,Chetyrkin:1981qh,Laporta:2000dsw} and generalized unitarity \cite{Bern:1994zx,Bern:1995db,Britto:2004nc,Britto:2006sj,Anastasiou:2006gt,Ossola:2006us}.
%
An efficient method to compute directly the 2-point function coefficients, projecting it out of a double-cut, relies on Stokes' theorem \cite{Mastrolia:2009dr}, and it is based on a reparameterization of the virtual spinors in terms of the external ones that is implemented via the following spinor rotation matrix:
\begin{eqnarray}
    \begin{pmatrix}
 \lambda_x \\
 \lambda_y
    \end{pmatrix}
     =  
   \frac{1}{\sqrt{1+ z \bar{z}}}\left(
\begin{array}{cc}
 1 & \bar{z} \\
 -z & 1 \\
\end{array}
\right)
\begin{pmatrix}
 \lambda_{a} \\
 \lambda_{b}
    \end{pmatrix}\,,
    \label{eq:stokes_par}
\end{eqnarray}
where $z$ and $\bar{z}$ are complex conjugate variables.
%
The integration measure is defined as:
\begin{eqnarray}
    \int d \text{LIPS}_2 = 
    \frac{-1}{8 \pi} 
    \oint \frac{dz}{2\pi i } 
    \int    
    \frac{d{\bar z}}{(1 + z \bar{z})^2} \, . 
\end{eqnarray}

The integrand is a generic rational function $g(z,\zb)$, which can be integrated, using complex analysis, by a two-step procedure: finding a primitive function in $\zb$, say $G(z,{\bar z})$ (considering $z$ as an independent variable), and applying Cauchy's residue theorem in $z$.
 The computation of the primitive generates two kinds of contributions: a rational part and a logarithmic one, i.e., 
 $G = G_{\rm rat} + G_{\log}$. 
 For the determination of the anomalous dimensions,
 it is sufficient to determine the UV singularities carried out by the 2-point integrals, neglecting the higher-point functions that contain only IR singularities.
 Since the cut of the 2-point scalar function is rational, the contribution of the 2-point integral within the double-cut of a generic one-loop integral appears exclusively in the rational contribution of the double-cut \cite{Mastrolia:2009dr}. 
 Therefore, in order to isolate the coefficient of the 2-point function, it is sufficient to retain only the rational part (Rat) of the phase-space integration, coming just from $G_{\rm rat}$, yielding:
    \begin{equation}
        {\rm Rat}\int d{\rm LIPS}_2 \ g(z,\zb)
        \equiv 
         \frac{-1}{8 \pi} 
         \oint \frac{dz }{ 2 \pi i} 
         G_{\rm rat}(z,{\bar z}) 
         =
         \frac{-1}{8 \pi} 
         \sum_{z_0\in {\cal P}_{G_{\rm rat}}} 
         \!\!\!
         \Res_{(z,\zb)=(z_0,z_0^*)} G_{\rm rat}(z,\zb)\, . 
    \end{equation}
  
The above formula considers the contributions of the residues at all the finite poles $z_0$ of $G_{\rm rat}$, collected in the set ${\cal P}_{G_{\rm rat}}$  
(when taking the residues, the variable ${\bar z}$  is evaluated at $z_0^* = (z_0)^{\rm c.c.}$). 
Stokes' integration method offers the advantage of projecting directly the contribution of the 2-point functions out of double-cuts, 
hence its application appears to be simpler than other techniques. 

Within the coefficient of the 2-point integrals, both UV and IR collinear divergences are combined.
Therefore, in order to isolate the genuine UV anomalous dimension,
one has to compute and add,
out of the complete IR anomalous dimension, only the collinear part.
